\section{Finite edge equations and the skew-form defect of a mask}
\label{sec:native-edge-isotropy}

The preceding abstract obstructions do not exhaust the structure of native
normal equations.  This section derives an additional identity from the
actual labelled face pairings, including the conventions needed to compare
it with the maintained cycle matrix.  The interior-edge formulation is
standard; see Burton's Definition~2.4 and finite-link lifting argument
\cite{edge-burton}.  The cancellation underlying the skew identity is the
local one in Neumann's proof that $\beta^*\beta=0$ \cite{edge-neumann}.
We use only that local cancellation, and check the finite-boundary case
explicitly.  No ideal-triangulation homology theorem is imported.

\subsection{The exact native incidence presentation}

Let $T$ be accepted by \path{normal_surface_geometry._prepare}.  It is a
connected, orientable finite source with reciprocal face pairings, circle
or interval edge links, sphere or disc vertex links, and one torus boundary.
Distinct faces of one tetrahedron may be paired, and local edges and vertices
may be identified; edge reversals and self-pairing of a single face are
excluded.  The argument does not require a single source vertex.

The native quadrilateral indices are
\[
 0:\{01,23\},\qquad 1:\{02,13\},\qquad 2:\{03,12\}.
\]
Reconstruct tetrahedron signs $\varepsilon_\tau\in\{1,-1\}$ by choosing
$\varepsilon_0=1$ and imposing
$\varepsilon_\sigma=-\operatorname{sgn}(p)\varepsilon_\tau$ across a face
pairing with vertex permutation $p$.  This is the orientation propagation
already validated by \path{_prepare}.  Its signs are not returned in the
prepared dictionary.  The returned \path{edge_orientations} array instead
records edge union-find parity and cannot replace them.

Let $I$ be the global edges absent from \path{boundary_incidence}.  Define
$G_{e,3\tau+k}$ to count local edges of tetrahedron $\tau$ belonging to the
global interior edge $e\in I$ and opposite-pair type $k$.  Count occurrences:
if both opposite edges have the same global image, their count is two.
Put
\[
 J=\begin{pmatrix}0&1&-1\\-1&0&1\\1&-1&0\end{pmatrix},\qquad
 \Omega=\bigoplus_\tau\varepsilon_\tau J,\qquad E=G\Omega.
\]
For $g_k=G_{e,3\tau+k}$ the local row is
\[
 E_{e,\tau}=\varepsilon_\tau(g_2-g_1,\ g_0-g_2,\ g_1-g_0).
\]
Equivalently, one occurrence of type $0$, $1$ or $2$ contributes
$\varepsilon_\tau(0,1,-1)$, $\varepsilon_\tau(-1,0,1)$ or
$\varepsilon_\tau(1,-1,0)$ respectively.  Thus
\[
 |E_{ej}|\le2,\qquad \sum_e|E_{ej}|\le4,
 \qquad \sum_{e,j}|E_{ej}|\le12t.
\]
Each Q type meets four local edges, so aggregation can only decrease these
bounds.  Every local triple of columns sums to zero.  This is a sparse
constant-coefficient presentation even when a cycle-forest presentation is
dense.

\begin{proposition}[Equality with the maintained matching obstruction]
Let $C$ be the full Q-cycle matrix constructed by either the native producer
or the independent original-STANDARD-row verifier.  Then
\[
 \ker E=\ker C,\qquad
 \operatorname{row}_{\mathbb Q}E=\operatorname{row}_{\mathbb Q}C.
\]
The row-space equality remains true after restriction to any Q column mask.
\end{proposition}
\begin{proof}
Across a paired face $(\tau,f)\to(\sigma,g)$, with $w=p(v)$, the original
matching row is
\[
 x_{\tau,v}+q_{\tau,\operatorname{quad}(f,v)}
   -x_{\sigma,w}-q_{\sigma,\operatorname{quad}(g,w)}=0.
\]
Its prescribed triangle-potential difference is exactly the link-dual edge
label used by both constructors.  Around a local edge $ab$, order the other
vertices $c,d$ so that $(a,b,c,d)$ has sign $\varepsilon_\tau$.
The local period contributes
$q_{\tau,\operatorname{quad}(a,c)}-q_{\tau,\operatorname{quad}(a,d)}$.
For $01,02,03$ this gives the three rows displayed above; the opposite
edges give the same rows.  Summing incidences gives the row of $E$.
More explicitly, exit through face $c$.  The next oriented edge-end tuple
is $(p(a),p(b),p(d),p(c))$, whose sign is the propagated tetrahedron sign.
Summing the original matching rows along this walk cancels every triangle
term and leaves precisely the displayed Q period, at either end of the
interior edge.

On a disc or sphere link, small loops around interior link vertices generate
the rational cycle space of its dual graph.  These vertices are the ends of
interior source edges.  Thus $Eq=0$ exactly when every prescribed period
vanishes, equivalently when signed triangle potentials exist.  Contracting
zero-labelled edges, taking fundamental cycles, primitive scaling and
deduplication preserve this period span.  Active self-loops retain their
labels, including when triangle terms cancel in an original matching row.
This proves equal kernels and rational row spaces.  Restricting the two
row-space transformations to any columns proves the last assertion.
\end{proof}

This is signed linear algebra; it imposes neither Q compatibility nor
positive Euler characteristic.  Sphere/disc links are essential: ideal
vertices can have additional homological periods.  Boundary edges have
interval stars and do not supply the closed-loop equations used here.

The equality need not be an integral unimodular row change.  If $E=UC$ and
an alternative LP presentation produced a negative dual $y$ with
$E^{\mathsf T}y\ge c$, the current cycle-row packet would require
$U^{\mathsf T}y$.  Using $y$ unchanged would alter its meaning.  Rational
elimination supplies polynomial-size row transformations, but this section
does not replace the maintained LP matrix or certificate schema.

\subsection{Full-source skew cancellation}

\begin{proposition}[Native full-source isotropy]
With the literal type grouping and orientation signs above,
\[
 G\Omega G^{\mathsf T}=0,\qquad
 E\Omega E^{\mathsf T}=0,\qquad
 C\Omega C^{\mathsf T}=0.
\]
\end{proposition}
\begin{proof}
Expand the first product over ordered local edge occurrences.  Equal or
opposite edges have equal types and contribute zero.  Adjacent edges share
one face.  A vertex permutation induces a permutation of the three
opposite-edge types with the same parity, so paired-face contributions are
opposite by
$\varepsilon_\sigma\operatorname{sgn}(p)=-\varepsilon_\tau$.
Unpaired boundary faces contain only boundary edges; they contribute no
term to the interior-edge rows.  This proves the first identity by pairing
local terms, before aggregating any repeated global identifications.

Now $\Omega^{\mathsf T}=-\Omega$ and $\Omega^3=-3\Omega$, so
\[
 E\Omega E^{\mathsf T}
   =-G\Omega^3G^{\mathsf T}=3G\Omega G^{\mathsf T}=0.
\]
Finally write $C=RE$ over $\mathbb Q$ and multiply the middle identity by
$R$ and $R^{\mathsf T}$.  Neither a forest choice nor primitive row
normalization affects the conclusion.
\end{proof}

The proof allows selfattachments between distinct faces and repeated edge
occurrences.  Including boundary-edge rows would leave boundary terms and
is not justified.  Nor can one arbitrarily regroup coordinates from
different tetrahedra while keeping the same form $\Omega$.

At a one-vertex, two-boundary-face stage, there are $t-1$ interior edges.
The finite matching-dimension calculation in
Section~\ref{sec:euler-transversal} gives $\operatorname{rank}E=t-1$ and
$\dim\ker C=2t+1$.  Quotienting by the $t$ block-ones directions leaves
dimension $t+1$, not a polylogarithmic space.  On the block-sum-zero
representatives, $\Omega^2=-3I$.  A row covector $r$ corresponds to
$a=\Omega r^{\mathsf T}/3$, for which $a^{\mathsf T}\Omega=r$.
The quotient kernel is the symplectic orthogonal complement of the
converted $(t-1)$-dimensional isotropic row space.  This conversion is
needed: an ordinary dot annihilator is not directly the skew annihilator
of the unconverted rows.  None of these dimensions bounds the number of
compatible positive rays.

\subsection{Two-type faces and the defect of further restriction}

Write $K_\tau=C_\tau\Omega_\tau C_\tau^{\mathsf T}$.  The sum of these
matrices is zero, but the individual summands need not vanish.  There is a
useful exact rule for every local mask:
\[
 C_{\tau,M}\Omega_{\tau,M M}C_{\tau,M}^{\mathsf T}=w_\tau K_\tau,
 \qquad
 w_\tau=\begin{cases}
 1&\text{three retained types},\\
 1/3&\text{two retained types},\\
 0&\text{at most one retained type}.
 \end{cases}
\]
Indeed, write $C_\tau=[a,b,c]$ with $c=-a-b$.  Direct multiplication gives
$K_\tau=3\varepsilon_\tau(ab^{\mathsf T}-ba^{\mathsf T})$.
Each of the three two-column principal restrictions gives exactly one
third of this matrix, independently of the omitted type.

\begin{proposition}[Every two-type face is isotropic]
If a mask retains exactly two types in every tetrahedron, its restricted
matrix is isotropic for the principal restricted form and has the same
rank as the full matrix.  At a one-vertex, two-boundary-face stage its
signed kernel has dimension $t+1$.  Restricting $s$ blocks of this mask
further to at most one type gives defect
\[
 -\frac13\sum_{\tau\in S}K_\tau,
 \qquad\text{of rank at most }2s.
\]
\end{proposition}
\begin{proof}
The full two-type product is $\frac13\sum_\tau K_\tau=0$.  Each omitted
column is minus the sum of its two retained columns, so the column span,
and hence the rank, is unchanged.  The stated nullity follows from
$2t-(t-1)=t+1$.  A further singleton or empty block has zero principal
skew product.  Subtraction gives the defect formula; every $K_\tau$ has
rank at most two.
\end{proof}

The automatic block-minimum normalization from
Section~\ref{sec:euler-transversal} motivates these faces, but there are
$3^t$ choices of omitted types.  A two-type face still permits incompatible
positive pairs.  The dimension and isotropy statements neither select
polynomially many faces nor resolve compatibility with the canonical Euler
objective.

\begin{proposition}[Mask defect]
Fix the full source matrix $C$ and a mask $M$.  If $b$ tetrahedral blocks
have been changed from their full three columns to proper subsets, then
\[
 \operatorname{rank}\bigl(C[:,M]\Omega[M,M]C[:,M]^{\mathsf T}\bigr)\le2b.
\]
In particular, for a set $S$ of singleton-restricted tetrahedra with all
other blocks full, the exact defect is
\[
 C[:,M]\Omega[M,M]C[:,M]^{\mathsf T}
       =-\sum_{\tau\in S}C_\tau\Omega_\tau C_\tau^{\mathsf T}.
\]
It need not vanish.
\end{proposition}
\begin{proof}
The exact weighted rule gives the general formula
$-\sum_\tau(1-w_\tau)K_\tau$.  A singleton's weight is zero, giving the
displayed special case.  For each proper subset the difference has image
contained in the column span of $C_\tau$, of dimension at most two.
Rank subadditivity proves the bound.
\end{proof}

For an exact algebraic counterexample to automatic inheritance, take
\[
 r_1=(1,-1,0\mid1,-1,0),\qquad
 r_2=(0,1,-1\mid0,-1,1),\qquad \Omega=J\oplus J.
\]
The full pairing is $3-3=0$.  Retaining the first block and only one
coordinate of the second leaves pairing $3$.  This is an abstract control,
not a native triangulation.  Thus uniform two-type restriction preserves
isotropy, whereas arbitrary mixtures of full, two-type and singleton
blocks require the weighted defect formula or another invariant.

\subsection{What the identity excludes and what remains open}

The periodic-grid matrix of Section~\ref{sec:sparse-coordinate-separators}
fails full-source isotropy for every choice of signs on its designated
triples.  Choose two horizontally adjacent grid vertices whose incidence
rows were not deleted.  They occur together in exactly one designated
triangle, where their row pieces are $(-1,0,1)$ and $(1,-1,0)$.
Their local skew pairing is $3$, so the full pairing is $\pm3$ with no
other block available to cancel it.  Changing the row basis cannot turn a
nonzero skew product into zero.  Thus this particular labelled kernel is
not a native full Q kernel with the same triple grouping.  This assertion
concerns the represented kernel, not every representation of an abstract
column matroid.

There is also a sparse-row exclusion for the earlier dense moment-curve
family at one-vertex stages.  Its nonzero row functionals have support at
least $3t-\lfloor3(t-2)/2\rfloor$: if $d$ full triples and $a$ other
columns vanish, independence forces $2d+a\le t-2$, hence at most
$3d+a\le\lfloor3(t-2)/2\rfloor$ zeros.  In contrast, the $t-1$ independent
native edge rows have total support at most $12t$, so one has support at
most $\lfloor12t/(t-1)\rfloor$.  These inequalities exclude that dense
family for $t\ge8$.  Neither abstract obstruction had claimed a native
realization; these additional constraints explain concretely why one
cannot infer such a realization from their earlier algebraic properties.

The useful open obligation is now sharper: obtain a search theorem from
the actual finite-link realization, its block-skew structure and the
canonical Euler objective, with valid control on descendant masks.
Sparsity, isotropy or the defect estimate by themselves do not give a
separator bound, a small resolving parameter or a general quasi-polynomial
recognizer.  No new pivot, source simplification, LP shortcut or runtime
claim is introduced here.

